\chapter{Conclusion and Future Work}

\section{Conclusion}

This project set out to build a fraud detection system for UPI transactions that goes beyond simple rules and beyond the standard approach of training a classifier on raw transaction features. Looking back at what we built, we are genuinely satisfied with how it came together.

The central thesis of An Intelligent Hybrid Machine Learning Framework for UPI Fraud Detection is that fraud detection works better when you think about it as a problem of understanding behavior, not just classifying individual data points. A single transaction is a limited view it tells you the amount, the time, the parties involved. But when you look at that transaction in context how does this amount compare to what this sender normally does? How many transactions has this sender made in the past hour? What does this sender's position in the broader transaction network look like? a much richer picture emerges, and that richness translates into better predictions.

The results back this up. Our final system achieves a ROC-AUC of approximately 0.97 and a fraud recall of over 91\%, which means it correctly identifies nine out of ten fraud cases. The ablation study showed that each component we added graph features, velocity modeling, behavioral profiling, Z-score anomaly detection, and cost-sensitive threshold optimization contributed meaningfully to the final result. Removing any one of them makes the system measurably worse.

The cost-sensitive threshold optimization deserves particular mention because it is the kind of insight that is easy to overlook if you are thinking purely about model accuracy. By explicitly modeling the financial cost of each type of error and choosing the threshold that minimizes total expected loss rather than maximizing F1-score, we reduced the projected financial loss on our test set by more than 50\% compared to using the same model with a default 0.5 threshold. This is a practical contribution that carries over directly to real-world deployment.

The deployment layer FastAPI for programmatic access and Streamlit for analyst interaction demonstrates that this is not just an academic pipeline. Both interfaces are working, tested, and reasonably production-ready. The API's average latency of 18 milliseconds is well within the time budget for inline UPI fraud screening.

In short, this project shows that with careful feature engineering, thoughtful handling of class imbalance, and explicit cost-aware decision making, it is possible to build a UPI fraud detector that substantially outperforms both rule-based baselines and simple ML approaches.

\section{Limitations}

Being honest about what the system does not do well is as important as documenting what it does well.

\begin{itemize}
    \item \textbf{Graph leakage:} The graph features are computed over the full dataset before splitting. This means that some information from test transactions technically leaks into the training features via the graph structure. In a true production setting, the graph would be built from historical data only, not from future transactions.

    \item \textbf{Static model:} The model is trained once and does not update as new transactions arrive. Fraud patterns evolve over time, so a model trained today will gradually become less effective as fraudsters adapt. Periodic retraining is necessary in production.


    \item \textbf{Single-hour velocity:} We compute velocity only in a one-hour window. Fraud patterns at different time scales (five minutes, one day, one week) might provide additional signal.

    \item \textbf{No streaming support:} The system processes transactions in batch (for training) or one-at-a-time (for API inference). True streaming fraud detection processing a continuous flow of transactions with real-time feature computation requires different infrastructure.
\end{itemize}

\section{Future Work}

Several directions stand out as natural extensions of this work.

\subsection{Graph Neural Networks (GNN)}

The graph features we currently compute (PageRank, degree centrality) are relatively simple network statistics. Graph Neural Networks, which learn representations of nodes by aggregating information from their neighborhood in the graph, can capture much more complex structural patterns. Papers like Hu et al. (2022) \cite{hu2022} have shown strong results applying GNNs to financial fraud detection. Replacing the hand-engineered graph features with learned GNN embeddings is a natural next step.

\subsection{Real-Time Streaming Detection}

The current system operates in batch mode for training and single-request mode for inference. A production fraud detection system would need to process a continuous stream of transactions, updating the velocity and behavioral features in real time as new transactions arrive. Apache Kafka for streaming infrastructure combined with a lightweight online learning model would be an appropriate technical approach.

\subsection{SHAP Explainability}

The system already includes a placeholder for SHAP (SHapley Additive exPlanations) integration in \texttt{src/explain.py}. SHAP provides model-agnostic feature attribution for any given prediction, it tells you exactly how much each feature contributed to pushing the probability up or down. This is valuable for two reasons: it helps analysts understand why a transaction was flagged (improving trust and compliance), and it helps identify when the model's reasoning is based on spurious correlations. Completing the SHAP integration is a priority for the next version.

\subsection{Multi-Algorithm Ensemble}

The current system uses XGBoost as the sole classifier. An ensemble that combines XGBoost with LightGBM and a shallow neural network via stacked generalization (as explored by Kamble et al.) would likely provide small but consistent improvements in precision and recall.

\subsection{Concept Drift Detection and Adaptive Retraining}

A production deployment needs mechanisms to detect when the model's performance is degrading over time (concept drift) and trigger retraining when it does. Statistical tests like Population Stability Index (PSI) on input feature distributions can serve as early warning systems for drift.

\subsection{Docker Deployment}

Currently, the system requires Python and all dependencies to be installed manually. Packaging the API server and dashboard as Docker containers would make deployment dramatically simpler and more reproducible. A \texttt{docker-compose.yml} file could bring up the entire system (API + dashboard) with a single command.

\subsection{Feedback Loop Integration}

In a real deployment, fraud analysts review flagged transactions and provide ground truth labels (confirmed fraud or false alarm). Feeding this analyst feedback back into the training data as new labeled examples would allow the system to continuously improve its understanding of current fraud patterns.

\section{Summary}

An Intelligent Hybrid Machine Learning Framework for UPI Fraud Detection demonstrates that a relatively small team, working with open-source tools and a publicly available dataset, can build a fraud detection system that achieves competitive performance through thoughtful engineering. The key lessons from this project are:

\begin{enumerate}
    \item Feature engineering matters more than algorithm selection. The graph features, velocity computation, and behavioral profiling together account for a larger fraction of the performance gain than switching from Random Forest to XGBoost.
    
    \item Cost-sensitive threshold optimization is not optional in fraud detection. Optimizing for accuracy or F1-score leaves significant financial value on the table.
    
    \item End-to-end thinking pays off. A system that produces great evaluation numbers but cannot be deployed is not useful. Investing in the API and dashboard layers made the work real.
    
    \item Honest evaluation matters. Running the final evaluation on a completely held-out test set and being transparent about the graph leakage limitation reflects the kind of scientific rigor that makes results trustworthy.
\end{enumerate}

We hope that this work can serve as a useful reference for others working on similar fraud detection problems in the Indian digital payments ecosystem.
